\newcommand{\NTracts}{9,794}
\newcommand{\NHU}{17.1}
\newcommand{\NPop}{41.8}
\newcommand{\NBldg}{19.2}
\newcommand{\NPoi}{2.0}
\newcommand{\RefCorr}{0.975}
\newcommand{\RefMedRatio}{1.05}
\newcommand{\RefGapShare}{1.9\%}
\newcommand{\RefGapRural}{4.7\%}
\newcommand{\RefGapUrban}{1.3\%}
\newcommand{\RefGapTribal}{4.3\%}
\newcommand{\RefGapNonTribal}{1.6\%}
\newcommand{\CovBldg}{90.0\%}
\newcommand{\CovPoi}{90.0\%}
\newcommand{\FlagBldg}{510}
\newcommand{\FlagPoi}{529}
\newcommand{\RTwoBldg}{0.93}
\newcommand{\RTwoPoi}{0.82}
\newcommand{\CovSpatRegMinBldg}{88.9\%}
\newcommand{\CovSpatRegMaxBldg}{93.2\%}
\newcommand{\CovMondRegMinBldg}{88.9\%}
\newcommand{\CovMondRegMaxBldg}{90.2\%}
\newcommand{\CovLoroBldg}{84.6\%}
\newcommand{\CovLoroMinBldg}{78.1\%}
\newcommand{\PrecBldg}{0.09}
\newcommand{\RecBldg}{0.26}
\newcommand{\PrevBldg}{1.9\%}
\newcommand{\LiftBldg}{5.1}
\newcommand{\PRAUCBldg}{0.10}
\newcommand{\PRAUCOlsBldg}{0.06}
\newcommand{\PRAUCRatioBldg}{0.08}
\newcommand{\PrecPoi}{0.58}
\newcommand{\PrevPoi}{24.0\%}
\newcommand{\DelThirtyBldg}{66\%}
\newcommand{\DelFiftyBldg}{96\%}
\newcommand{\DelThirtyPoi}{18\%}
\newcommand{\DelFiftyPoi}{56\%}
\newcommand{\CountyRho}{0.85}
\newcommand{\CountyN}{324}
\newcommand{\NFdrFamily}{25}
\newcommand{\PoiSviHi}{3.8\%}
\newcommand{\PoiSviLo}{8.9\%}
\newcommand{\PoiSviAor}{0.41}
\newcommand{\HFiveWildRural}{4.8}
\newcommand{\HFiveWildRuralCI}{[2.4, 9.4]}
\newcommand{\HFiveDroughtRural}{2.9}
\newcommand{\HFiveDroughtRuralCI}{[1.2, 7.1]}
\newcommand{\HFiveHeatTribal}{0.08}
\newcommand{\HFiveHeatTribalCI}{[0.0, 0.4]}
\newcommand{\HEAll}{44,945}
\newcommand{\HFourDprBldg}{-0.010}
\newcommand{\HFourDprBldgCI}{[-0.035, +0.016]}
\newcommand{\HFourDaucPoi}{+0.004}
\newcommand{\HFourDaucPoiCI}{[+0.001, +0.007]}
\newcommand{\GapMsft}{2.8\%}
\newcommand{\GapOsm}{1.3\%}
\newcommand{\GainBldg}{1.7}
\newcommand{\OsmShareA}{31\%}
\newcommand{\OsmShareC}{51\%}
\newcommand{\GapShareA}{7.0\%}
\newcommand{\GapShareB}{1.8\%}
\newcommand{\GapRuralA}{15.7\%}
\newcommand{\GapRuralC}{4.2\%}
\newcommand{\GapSviC}{5.3\%}
\newcommand{\DidRuralPP}{9.9}
\newcommand{\DidRuralCI}{[-15.3, -5.5]}
\newcommand{\DidRuralP}{0.003}
\newcommand{\DidTribalPP}{6.7}
\newcommand{\DidTribalCI}{[-14.3, -0.6]}
\newcommand{\DidTribalP}{0.061}
\newcommand{\HThreeRural}{$+$0.060 log points, 95\% CI 0.013--0.106}
\newcommand{\HThreeRuralP}{0.040}
\newcommand{\EOneLost}{179,000}
\newcommand{\EOneTracts}{186}
\newcommand{\EOneBexar}{181}
\newcommand{\EOneResidents}{889,000}
\newcommand{\EOneSviShare}{63\%}
\newcommand{\EOneSviAll}{34\%}
\newcommand{\EOneFlagAfter}{82\%}
\newcommand{\EOneFlagBefore}{2.7\%}
\newcommand{\EOneFlagOther}{4.5\%}
\newcommand{\EOneRecall}{83\%}
\newcommand{\EOneRefGapAfter}{99\%}
\newcommand{\EThreeSummary}{rural +0.034 (endpoint 2026-08-19; 95\% CI 0.003--0.065); rural +0.040 (excluding Bexar County; 95\% CI 0.008--0.072); tribal +0.035 (endpoint 2026-08-19; 95\% CI 0.004--0.067); tribal +0.036 (excluding Bexar County; 95\% CI 0.002--0.070)}
\newcommand{\ETwoText}{For rural vs.\ urban tracts the adjusted odds ratio of a reference gap is 3.86 [2.40, 6.22], but 1.06 [0.72, 1.57] for the map-trained flag and 0.83 [0.48, 1.44] for the reference-anchored flag; for tribal vs.\ non-tribal tracts the values are 4.17 [1.89, 9.23], 1.26 [0.38, 4.20] and 0.92 [0.17, 5.14]. \emph{Absorption} is visible in the expectation itself: relative to the reference, observed completeness in rural tracts is lower by -0.031 [-0.053, -0.009] log points (region FE), and the map-trained expectation is already lower by -0.025 [-0.053, 0.004], so it reproduces 80\% of the rural shortfall (64\% of the tribal one). \emph{Resolution} limits even the reference-anchored model: its tract-level residual SD is 0.21 log points (map-trained: 0.19), against 0.22 for a 20\% gap, and it flags only 2.3\% of tracts (precision 0.16, recall 0.20 for reference gaps; map-trained 0.09 and 0.26).}
\newcommand{\OrigHOne}{4.9 vs.\ 2.3}
\newcommand{\OrigHOneP}{0.0004}
\newcommand{\RobustText}{Across 7 specifications (primary; no housing-type covariates; negative-binomial GLM; 80\% and 95\% ranges; within-region calibration; excluding tracts with fewer than 50 housing units), 53 of 56 flag contrasts keep the significance status of the primary analysis. The lower POI flag rate in high-SVI tracts holds in every specification (risk difference -0.040 to -0.067). The exceptions are: no housing-type covariates, building svi: RD +0.059 [+0.025, +0.086]; negative-binomial GLM, POI rural: RD -0.010 [-0.024, +0.002]; 95\% range, POI tribal: RD -0.012 [-0.028, -0.002]. Most notably, dropping units-in-structure covariates makes building flags significantly more common in high-SVI tracts (aOR 2.91), consistent with absorption: mobile-home and multi-unit housing shares, which correlate with SVI, explain away part of the high-SVI shortfall in the primary model. Details are in Appendix Table~\ref{tab:robust}.}
\newcommand{\RobustTextShort}{Across 7 specifications (primary; no housing-type covariates; negative-binomial GLM; 80\% and 95\% ranges; within-region calibration; excluding tracts with fewer than 50 housing units), 53 of 56 flag contrasts keep the significance status of the primary analysis. The lower POI flag rate in high-SVI tracts holds in every specification (risk difference -0.040 to -0.067). The exceptions are: no housing-type covariates, building svi: RD +0.059 [+0.025, +0.086]; negative-binomial GLM, POI rural: RD -0.010 [-0.024, +0.002]; 95\% range, POI tribal: RD -0.012 [-0.028, -0.002]. Most notably, dropping units-in-structure covariates makes building flags significantly more common in high-SVI tracts (aOR 2.91), consistent with absorption: mobile-home and multi-unit housing shares, which correlate with SVI, explain away part of the high-SVI shortfall in the primary model. Full results are in \texttt{outputs/tables/robustness.csv}.}
